% Some local defs...
\def\struct{{\tt struct}}
\def\LAYER{{\tt LAYER}}
\def\NEURON{{\tt NEURON}}
\def\WEIGHT{{\tt WEIGHT}}

\chap Software Construction

In this chapter I will outline the construction of the code and the
data structures used to implement the back-propagation algorithm.
The package has been written entirely in C and all references to C
``objects'' such as keywords or variable names
will be in {\tt typewriter} style type.
All developmental work was carried out on an IBM PC, using Microsoft
QuickC$^{\rm TM}$.
The software has been sucessfully ported without modification to a
Pyramid super-mini running Unix Version V.
Execution speeds will be discussed in a later chapter, entitled
``Performance''.

This chapter is divided into two parts.
The first part contains a description of the most important design
decision, namely the data structures used for the dynamically
allocated network.
The second section details a rough outline of the flow of control of
the software, especially as an aid to someone wishing to enhance the
code.
This section also details some important features of the code.

A printout of the C source code appears in toto in Appendix
A.
The source, include, object and executable files (as well as the
\TeX\ source for this document) are also on a diskette inside the
back cover of this work.

\sec Data Structure Design

As stated in the software specification, the code is to
operate with networks of variable size and configuration, restricted
only by available memory of the particular implementation.
Because of this constraint, all allocation of space for the network
(\neus\ and connections) is completely dynamic.
The final result is a dynamically allocated structure that has no
limits of size except that it must have fewer than~$8$ layers.
This is not an unreasonable constraint, as even during testing
no more than~$4$ layer networks were ever used, and some authors suggest
that there need
never be more than~$4$ layers for {\sl any} problem.

In all sections following, the makeup of the individual C structures
will be discussed, with field meanings, names and types given.
The format followed will be --- structure type name: field$_0$
description, type, name; field$_1$ description, type, name; etc.
First I will describe the most important structure, namely that used
for the network, then the other structures used for different
purposes within the network.

\ssec Network Structure

A back-propagation network can be divided into~$3$ parts.
A network is made up of layers.
Each layer contains \neus.
And each \neu\ has connections (weights) to the layer above.
The most important constraint to the problem of a data structure for
a network is the fact that it must be variable in size for each run
(this also allows the software to deal with networks as large as
available memory in any given system).
Because of this the structure is made up of dynamically allocated
arrays of pointers to pointers to structures.

\sssec Weights

Weights are stored in a \struct\ of type \WEIGHT.
Contained in \WEIGHT\ are fields for: weight or connection
strength, {\tt double w}; $\partial E \over \partial w$, {\tt double
dE\_dw}; and $\Delta w(t-1)$, {\tt double old\_delta\_w}.

\sssec Neurons

Neurons are stored in a \struct\ of type \NEURON.
\NEURON\ contains fields for: output of the neuron, {\tt double
y}; a pointer to an array of weights for this neuron's connection 
to the layer above, {\tt WEIGHT *w}.

\sssec Layer

Layers are stored in a \struct\ of type \LAYER.
A \LAYER\ is made up of fields for: a pointer to an array of \neus\
for that level, {\tt NEURON *neuron}.

\sssec Network

A network is stored as a pointer to an array of type \LAYER\ (that
is, {\tt LAYER *}).
Throughout the code, this is how references to the network are
passed.

\ssec Structures Used in Back-Propagation Calculations

Large numbers of intermediate variables are necessary during
calculations using the back-propagation method.
All are stored as dynamic arrays.

\sssec Temporary Values of $\partial E \over \partial x$

Temporary values of $\partial E \over \partial x$ are stored in an
array of type {\tt double}.

\sssec Temporary Values of $\partial E \over \partial y$

Temporary values of $\partial E \over \partial y$ are stored in an
array of type {\tt double}.

\ssec Input/Output Data Structures

To minimize relatively slow disk input/output for such a repetitive
action as learning, all input/output data for learning and all input
data for operation is read in at once when it is needed.

\sssec Neuron's Input/Output Value

Every \neu's input value for layer~$0$ \neus\ and every \neu's
expected
output for top layer \neus\ is stored in \struct s of type {\tt
OP\_DATA}.
Contained in {\tt OP\_DATA} is a field {\tt double x}, which is the
value of input or output for the \neu\ in question.
For each I/O case there is an array of {\tt OP\_DATA}.

\sssec Input/Output Cases

Input and output cases are stored in a \struct\ called {\tt I\_O}.
This structure contains a pointer to a particular case, {\tt
OP\_DATA *item}.
All I/O data is stored in arrays of {\tt I\_O}, each element of
which contains a pointer to a particular case.

\ssec General Data Structures

\sssec Data File Information

Every data file used by fishNET has its information stored in a
\struct\ of type {\tt DATA\_FILE}.
Contained in {\tt DATA\_FILE} are: the file name, {\tt char name[35]};
a single letter abbreviation of the file's type, {\tt char type};
the number of input or output cases for this file, {\tt int
i\_o\_cases}; and the number of \neus\ in the layer this file is
associated with, {\tt int \neus}.

\sssec Most Commonly Used Parameters

During the operation of fishNET, all constantly used data is stored
in one \struct, called {\tt QUESTION}.
Throughout the program, pointers to {\tt QUESTION} ({\tt QUESTION *}) 
are used to pass important data.
{\tt QUESTION} contains: the number of layers in the network,
{\tt int layers}; the number of \neus\ in every layer, {\tt int
per\_layer[7]}; the nominal output width, {\tt int out\_width};
the maximum number of learning sweeps, {\tt int max\_sweeps};
the learning parameter $\alpha$, {\tt double alpha};
the learning parameter $\varepsilon$, {\tt double epsilon};
the name of the file where the network is to be saved, {\tt char
save\_name[35]};
a comment for the saved file, {\tt char save\_comment[79]};
a chunk of information about network teaching and execution files,
{\tt DATA\_FILE sample\_in};
{\tt DATA\_FILE sample\_out};
{\tt DATA\_FILE data\_in};
{\tt DATA\_FILE data\_out}.

\sssec Structure for Network Files

When a network file is loaded, it contains the complete network
(including weights), plus all the parameters necessary for
operation.
The function which reads in the network files returns a pointer to a
struct called {\tt LOAD\_BOTH}.
Contained in {\tt LOAD\_BOTH} are: 
a pointer to a network, {\tt LAYER *network};
a pointer to a structure containing all run time information, {\tt
QUESTION *parameter};
and the value of $t$ at which the network will start teaching, {\tt
int start\_time}.

\sec Program Design

There are two important functions which the program performs, and
they are: training artificial \nn s, and running \nn s.
The code is written so that both these functions are easily and
possibly separately performed.
Once the program has been given the necessary dimensions and
parameters for the network, it will teach it. 
Given the network
that you want executed, it will apply the required data to the
input, calculate, then save the output.
Hence, the program looks like:
$$
\matrix{{\rm expected\ data,}&&&&{\rm neural}\cr
	{\rm configuration\ data}&\mapright{}&{\rm teacher}
		&\mapright{}&{\rm `operator'}&\mapright{}
		&{\rm output}\cr
	&&\mapdown{}&&\mapup{}\cr
	&&{\rm network\ file}&&{\rm run\ data}\cr}
$$

The items configuration and expected data, network file, 
run data, and output have been discussed in the previous chapter.
Each of the other parts, the teacher and the neural `engine' will be
discussed below.

\ssec The Teacher

Teaching a network is broken into several phases.
For each phase, the appropriate subroutine and its functions will be
examined.

\sssec Reading the Parameters of Network Operation

The action taken here depends upon the command line options set.
If there are no options set, the program calls
$$
	{\tt parameter = input\_parameters();}
$$
which returns a pointer to a structure containing the important
parameters of network operation, read from the keyboard.

If the $-c$ option was used, the program instead calls
$$
	{\tt parameter = finput\_parameters(config\_file);}
$$
which returns a pointer to a structure, containing the important
parameters of network operation, read from the file whose name is
held in the variable {\tt char *config\_file}.

If the $-n$ option was used, fishNET loads the network at the same
time as the operational parameters, so please see the next section
on allocating space for the network.

\sssec Allocating Space for The Network

If the network is to be trained from the start (that is, no network
is pre-loaded), the program allocates space for it by
$$
	{\tt network = allocate(parameter);}
$$
This function allocates space (from the heap) and sets up random
weights for a network of the desired size, and returns a pointer to
it ({\tt LAYER *}).

If the network is to be loaded from disk (the $-n$ command line
directive was used), fishNET reads the parameters and allocates
space for the network by
$$
	{\tt net\_and\_param = load\_network\_and\_parameters(netfile);}
$$

The function takes as its arguments the name of the network file,
and returns a pointer to a structure which contains: a pointer to a
structure which holds the parameters; a pointer to a network, with
the values read from the file as the weights; and an integer value
of the time~$t$ the network is to start learning at.
For more information about the data structures used see 4.1.~Data 
structure design.

\sssec Loading the Expected Data

Once the network has been allocated, the data is read in and stored
in the structure previously described in 4.1.3.~Input/Output data
structures.
The function call returns a pointer to an array of type {\tt I\_O},
that is {\tt I\_O *}.
The call is
$$
\eqalign{{\tt input} &{\tt = get\_data(parameter->sample\_in);}\cr
	{\tt output} &{\tt = get\_data(parameter->sample\_out);}\cr
}$$

\sssec Teaching the Network

The network is taught by the function
$$
	{\tt learn(network,\ parameter,\ start\_time);}
$$
This function uses the information in the variable {\tt parameter},
which is actually a pointer to a structure of type {\tt QUESTION},
to modify the network, which is pointed to by {\tt LAYER *network}.
The initial value of~$t$ used by the routine is {\tt int
start\_time}.
The algorithms are explained in Chapter 2, and are implemented exactly
as documented there.

If the first algorithm is being used, ($-a$ specified, or the command
line default,) {\tt learn} calls the routines
{\tt operate}, {\tt back\_propagate}, and {\tt apply\_delta\_w} to
implement the algorithm.
Actually, {\tt apply\_delta\_w} is executed first, because a {\tt do
\dots while} loop is used to test the number of errors, which is not
known until the routine {\tt back\_propagate} is called.
Hence to make sure that {\tt apply\_delta\_w} is not called once too
often, (ie when the error is actually zero,) it is run first.
This works because as the network is allocated by fishNET, it is 
initialised and so the first running of {\tt apply\_delta\_w} does
nothing.

The routine {\tt operate} is detailed in the next part, The neural
`engine'.

{\tt Back\_propagate} calculates the error $\partial E \over
\partial w$ for each case and adds it to the total kept in memory
for every weight.
The error is calculated by comparing the output with the expected
output, and is described earlier.

{\tt Apply\_delta\_w} is used to calculate and then add $\Delta
w(t)$ for every weight in the network.

If the second algorithm is used (the $-1$ flag was used), the functions 
called are slightly different.
First, {\tt learn} calls {\tt operate}, then a function called
{\tt back\-\_propagate\-\_apply\-\_delta\-\_w}.
This routine is a combination of the two like--named routines
described above, {\tt back\-\_propagate} and {\tt apply\-\_delta\-\_w}.
It operates by calculating $\partial E \over \partial w$ for every
case and using this to calculate $\Delta w(t)$ and apply it for
every case.
This appears to make the network converge much more slowly, but is
included for experimental purposes.

\sssec Saving the Network

Once the network is fully `taught', or the maximum value of~$t$ is
reached, the action depends upon the setting of the command line
options $-d$, $-s$, and $-x$.

If $-d$ was specified, execution occurs immediately and the network
is discarded.

If $-s$ was selected, the network is saved to disk by a call to the
function
$$
	{\tt store\_network(network,\ parameter);}
$$
and execution occurs immediately.

If $-x$ was specified, the learning parameters are saved by
$$
	{\tt store\_learnt\_parameters(parameter);}
$$
and the program exits without execution.

\ssec The Neural `Engine'

\sssec Operating With a Network

The network passed to the function has the input applied to the
bottom layer and the output is stored in the top (output) layer of
the network.
The input of every \neu\ is the sum of the outputs of all the
\neus\ below times their interconnecting weights.
The output of every \neu\ is its input times the transfer function.
The call that performs these functions is
$$
	{\tt operate(network,\ input\_case->item,\ parameter);}
$$

